%% kapa: exam=1 round=24 year=2013 period=2 form=A subject=appraisal_law subject_ko=부동산관계법규 seq=25 number=65
%% source: Q-Net 감정평가사 기출문제 (cst00309) artlSeq=5073802 file=extracted/1008/5073802/2013년 제24회 감정평가사 1차시험 기출문제/2교시(A형).hwp
%% conversion: kordoc 4.24.51 -> tools/gen1.py (auto); review=proofread-sonnet; flags=none
\begin{question}{65}
국유재산법령상 시효취득에 관한 설명으로 옳은 것은? (다툼이 있으면 판례에 의함)
\begin{choices}
  \item 보존용재산은 「민법」상 시효취득의 대상이 된다.
  \item 일반재산은 취득시효기간 동안 계속하여 일반재산이어야만 취득시효가 완성된다.
  \item 행정재산이 기능을 상실하여 본래의 용도에 제공되지 않는 상태에 있었다면 공용폐지 없이도 시효취득의 대상이 된다.
  \item 토지의 지목이 도로이고 국유재산대장에 등재되어 있다면 그 사실만으로 당해 토지는 시효취득의 대상이 되지 않는다.
  \item 국유 하천부지는 공용개시행위가 있어야만 시효취득의 대상이 되지 않는다.
\end{choices}
\end{question}
\answer{2}
